\chapter{Examples and machine learning framework}
\label{sec:ch1-intro}

The object of this book is to present several machine learning procedures
providing at the same time context on the applicability domain so that the
reader will acquire an understanding of what algorithm is best for any given
practical situation.

We will start with the simplest one, namely the k-nearest-neighbors in
order to see where is the origin of the machine learning 'magic'; this will allow
to distinguish the distance based methods (usually relevant in low dimensions)
from protocols used in higher dimensions. Then we continue with ridge and
lasso regularization, naive Bayes estimators and neural networks; we close with
the K-means algorithm.

In order to situate our framework, the rest of this section is devoted to
some examples of applications. First we give a list and then make clear the
mathematical framework and the steps taken to address a machine learning
task.

\section{Examples of applications}
\label{sec:ch1-examples}

\subsection{Image Recognition and Computer Vision}
\label{sec:ch1-imgrec}

Machine learning has revolutionized the field of computer vision, enabling systems
to interpret and analyze visual data effectively. Several applications of
computer vision and their corresponding use cases are presented in Table~\ref{tab:ch1-cv}.
These include object detection and recognition, which identifies and classifies
objects in images or videos, enabling use cases such as autonomous driving
and retail inventory management through models like YOLO and Faster R-CNN.
Facial recognition is employed for identity verification in smartphones,
social media tagging, and law enforcement. In healthcare, medical image analysis
leverages convolutional neural networks to interpret scans (e.g., X-rays,
MRIs) for disease detection. Image segmentation supports pixel-level classification
for tasks such as road detection in autonomous vehicles and organ
delineation in medical imaging, using architectures like U-Net and Mask R-CNN.
Finally, super-resolution and image enhancement improve image quality
for applications ranging from satellite imaging to consumer photo restoration,
often utilizing generative models such as SRGAN.

\begin{table}[h]
\centering
\begin{tabular}{p{0.28\linewidth} p{0.62\linewidth}}
\toprule
\textbf{Application} & \textbf{Description and Use Cases} \\
\midrule
Object Detection and Recognition &
Identifies and classifies objects within images or videos. Used
in self-driving cars (detecting pedestrians, vehicles, traffic
lights), retail (inventory management), and leveraging models like YOLO and Faster R-CNN. \\
\midrule
Facial Recognition &
Identifies or verifies individuals based on facial features. Applications
include smartphone unlocking, social media photo
tagging (e.g., Facebook), and law enforcement for criminal
investigations. \\
\midrule
Medical Image Analysis &
Examines medical scans (X-rays, MRIs, CT scans) to detect
diseases or abnormalities. Assists in diagnosing conditions
like cancer using models like Convolutional Neural Networks
(CNNs). \\
\midrule
Image Segmentation &
Divides images into meaningful parts for pixel-level classification.
Used in autonomous vehicles (road and obstacle
detection) and healthcare (highlighting organs or brain tumors in scans) with algorithms like U-Net and Mask R-CNN. \\
\midrule
Super-Resolution and Image Enhancement &
Improves image quality by increasing resolution or restoring
damaged areas. Applications include satellite imaging and
consumer tools for photo enhancement using models like SRGAN. \\
\bottomrule
\end{tabular}
\caption{Applications of Machine Learning in Computer Vision.}
\label{tab:ch1-cv}
\end{table}

\subsection{Natural Language Processing (NLP)}
\label{sec:ch1-nlp}

Natural Language Processing (NLP) encompasses a range of techniques that
enable machines to understand, interpret, and generate human language. The
Table~\ref{tab:ch1-nlp} summarizes key applications of NLP and their corresponding use
cases. Text generation and language modeling automatically produce human-like
text from prompts using models such as GPT, powering tools like AI writing
assistants, email automation, and creative writing. Sentiment analysis classifies
text as positive, negative, or neutral, supporting social media monitoring,
product review evaluation, and customer feedback analysis through models like
BERT and LSTM. Machine translation enables multilingual communication by
translating text between languages, as seen in services like Google Translate,
leveraging Transformer architectures for improved accuracy. Named Entity
Recognition (NER) extracts and categorizes entities such as names, dates, and
locations, aiding information retrieval, legal document analysis, and healthcare
data processing. Finally, speech recognition and conversational AI convert
spoken language to text and generate responses, forming the backbone
of virtual assistants (e.g., Siri, Alexa) and customer service chatbots, using
advanced models like Transformers for text understanding and WaveNet for
speech synthesis.

\begin{table}[h]
\centering
\begin{tabular}{p{0.28\linewidth} p{0.62\linewidth}}
\toprule
\textbf{Application} & \textbf{Description and Use Cases} \\
\midrule
Text Generation and Language Modeling &
Automatically generates human-like text based on prompts,
using models like GPT. Applications include AI writing assistants
(e.g., ChatGPT), automatic email generation, and
creative writing. \\
\midrule
Sentiment Analysis &
Determines the sentiment (positive, negative, or neutral) in
text. Used in social media monitoring, product reviews, and
customer feedback analysis, leveraging models such as BERT
and LSTM. \\
\midrule
Machine Translation &
Translates text between languages, as seen in tools like
Google Translate. Transformer-based models enhance translation
accuracy, facilitating international communication,
global business, and tourism. \\
\midrule
Named Entity Recognition (NER) &
Identifies and classifies entities (e.g., names, dates, locations)
in text. Applications include information extraction,
legal text analysis, healthcare (e.g., extracting medical conditions),
and enhancing search engine capabilities. \\
\midrule
Speech Recognition and Conversational AI &
Converts spoken language to text (speech-to-text) and generates
meaningful responses (text-to-speech). Used in virtual
assistants (e.g., Siri, Alexa) and chatbots for customer support,
leveraging Transformer models for text understanding
and WaveNet for speech generation. \\
\bottomrule
\end{tabular}
\caption{Applications of Natural Language Processing (NLP).}
\label{tab:ch1-nlp}
\end{table}

\subsection{Generative AI}
\label{sec:ch1-genai}

Generative AI leverages machine learning models to create new content across
various domains, from images and text to music and 3D objects, see Table~\ref{tab:ch1-genai}
for some examples of use. Image generation creates realistic visuals or artistic
content from text or random noise using models such as GANs, powering
tools like DeepArt and DALL-E and enabling deepfake video synthesis. Text
generation produces coherent, context-aware text through language models
like GPT, supporting AI-driven content creation, writing assistants, code generation
(e.g., GitHub Copilot), and conversational agents. Music and audio
generation involves composing music or generating synthetic speech and sound
effects, with models like OpenAI's Jukebox and AIVA for music, and WaveNet
for high-quality text-to-speech in virtual assistants and audiobooks. Data augmentation
for training leverages synthetic data generation, often via GANs, to
improve model performance in scenarios with limited real-world data, particularly
in image recognition tasks. Lastly, 3D object and scene generation creates
virtual environments and models for gaming, simulations, VR, and AR, using
techniques such as 3D GANs to produce realistic objects and immersive worlds.

\begin{table}[h]
\centering
\begin{tabular}{p{0.28\linewidth} p{0.62\linewidth}}
\toprule
\textbf{Application} & \textbf{Description and Use Cases} \\
\midrule
Image Generation &
Creates realistic images or art from text or random noise
using models like GANs. Applications include tools like
DeepArt and DALL-E for art creation, as well as deepfake
video generation for realistic faces or scenes. \\
\midrule
Text Generation &
Produces coherent and contextually relevant text using language
models like GPT. Use cases include AI-driven content creation, writing assistants (e.g., Grammarly), automatic
code generation (e.g., GitHub Copilot), and chatbot
dialogues. \\
\midrule
Music and Audio Generation &
Composes music or generates synthetic speech and sound
effects. Models like OpenAI's Jukebox and AIVA are used
for creating music, while WaveNet powers high-quality text-to-speech
for virtual assistants and audiobooks. \\
\midrule
Data Augmentation for Training &
Generates synthetic data to enhance machine learning model
performance, especially in scenarios with limited real data.
GANs are used to create new data points, improving model
accuracy in tasks like image recognition. \\
\midrule
3D Object and Scene Generation &
Creates 3D models and virtual environments for gaming, simulations,
VR, and AR applications. Models like 3D GANs
are applied in video game design and movie special effects to
generate realistic objects and virtual worlds. \\
\bottomrule
\end{tabular}
\caption{Applications of Generative AI.}
\label{tab:ch1-genai}
\end{table}

\subsection{Examples: Game Play and Robot Training}
\label{sec:ch1-rl}

Machine learning and reinforcement learning (RL) have transformed gameplay
and robotics, enabling advanced training for AI agents and robots. Table~\ref{tab:ch1-rl}
contains a description of some key applications: Gameplay and AI for games
uses RL to achieve expert-level performance in titles like Chess, Go, and Dota
2 (e.g., AlphaGo, OpenAI Five) and to create adaptive NPC behavior. Robot
training and control applies RL for locomotion, object manipulation, and navigation
with tools like OpenAI Gymnasium and Mujoco, as seen in Boston
Dynamics' Spot. In autonomous vehicles, RL supports safe navigation and
decision-making for self-driving cars used by companies such as Waymo and
Tesla. Personalized game recommendations enhance engagement by adjusting
difficulty and suggesting challenges based on player behavior. Finally, sim-to-real
transfer enables robots to learn in simulations and apply skills in real-world
tasks, reducing cost and risk in areas like drone control and industrial
automation.

\begin{table}[h]
\centering
\begin{tabular}{p{0.28\linewidth} p{0.62\linewidth}}
\toprule
\textbf{Application} & \textbf{Description and Use Cases} \\
\midrule
Gameplay and AI for Games &
RL trains AI to play games like Chess, Go, and Dota 2 at
expert levels (e.g., AlphaGo, OpenAI Five). It is also used
for adaptive NPC behavior in game development. \\
\midrule
Robot Training and Control &
RL teaches robots tasks such as walking, object manipulation,
and navigation using tools like OpenAI Gymnasium
and Mujoco. Examples include Boston Dynamics' Spot. \\
\midrule
Autonomous Vehicles &
RL trains self-driving cars for safe navigation, obstacle avoidance,
and decision-making. Companies like Waymo and
Tesla rely on RL for vehicle optimization. \\
\midrule
Personalized Game Recommendations &
RL personalizes in-game experiences by adjusting difficulty
or suggesting challenges based on player behavior, enhancing
engagement in games like Minecraft and Fortnite. \\
\midrule
Sim-to-Real Transfer in Robotics &
RL enables robots to learn in simulations and apply learned
skills in real-world tasks, reducing risks and costs. Applications
include drone control and industrial robotics. \\
\bottomrule
\end{tabular}
\caption{Applications of RL in Games and Robotics.}
\label{tab:ch1-rl}
\end{table}

\subsection{Network-Based Recommendations and Relationships}
\label{sec:ch1-graph}

Graph-based techniques leverage the interconnected structure of data to provide
insights and improve applications cf. Table~\ref{tab:ch1-graph}. Recommender systems
use Graph Neural Networks (GNNs) to model user-item relationships, improving
recommendations on platforms like Netflix and Amazon. Social network
analysis applies GNNs to user connections for better ad targeting and content
delivery. In fraud detection, transaction networks modeled as graphs help
identify suspicious patterns for financial security. Finally, knowledge graphs
represent entities and relationships to enhance search engines such as Google,
enabling richer and more relevant results.

\begin{table}[h]
\centering
\begin{tabular}{p{0.28\linewidth} p{0.62\linewidth}}
\toprule
\textbf{Application} & \textbf{Description and Use Cases} \\
\midrule
Recommender Systems &
GNNs analyze user-item relationships to enhance recommendations
on platforms like Netflix, Spotify, and Amazon by
modeling interactions as graphs. \\
\midrule
Social Network Analysis &
GNNs analyze user connections on platforms like Facebook
and Twitter to improve ad targeting and content delivery by
modeling user interactions. \\
\midrule
Fraud Detection &
Transaction networks modeled as graphs help financial institutions
detect unusual patterns and relationships indicative
of fraudulent activities. \\
\midrule
Knowledge Graphs &
Graphs representing entities and their relationships enhance
search engines like Google, enabling richer and more relevant
search results. \\
\bottomrule
\end{tabular}
\caption{Applications of Graph-Based Techniques.}
\label{tab:ch1-graph}
\end{table}

\section{Summary: types of Machine Learning}
\label{sec:ch1-summary}

Summarizing the previous examples, machine learning can be broadly categorized
into three types (combinations or extensions are allowed):
\begin{itemize}
  \item \emph{supervised learning,} where models are trained on labeled data to learn mappings
  from inputs to outputs (e.g., predicting house prices from features);
  \item \emph{unsupervised learning,} where models find hidden patterns or structures in
  unlabeled data (e.g., clustering customers based on behavior); also in this denomination
  are generative models that create new objects similar to existing
  ones;
  \item \emph{reinforcement learning,} where agents learn to make sequential decisions by
  interacting with an environment to maximize a reward signal (e.g., training
  robots to navigate a maze).
\end{itemize}

\section{Supervised learning example: MNIST classification}
\label{sec:ch1-mnist}

The MNIST (Modified National Institute of Standards and Technology) dataset
is a widely used benchmark dataset in the field of machine learning and computer
vision, particularly for evaluating image classification algorithms \cite{turinici6}.
It consists of handwritten digits and serves as a fundamental dataset for developing
and testing various models; there are $N = 70{,}000$ images in the dataset,
which are usually split into two sets: $60{,}000$ training images and $10{,}000$ testing
images. Each image is a $28 \times 28$ matrix representing the pixel values of the
images, so the dimension of each data point is $d = 28 \times 28 = 784$; often this is
presented as flattened into a vector in $\R^{784}$.

Each image $X_i \in \R^{784}$ represented as a vector, contains the pixel intensity
values: $X_{ij} \in \{0, 1, \ldots, 255\}$ denotes the pixel value at the $j$-th position of
the $i$-th image. Often, the pixel values are divided by $255$ so that one obtains
values in $X_i \in [0,1]^{784}$. The image $X_i$ comes with a label $Y_i \in \{0, 1, 2, \ldots, 9\}$
representing the digit depicted in the image. Figure~\ref{fig:ch1-mnist} shows a selection
of $15$ images and their labels from the MNIST dataset arranged in a 3x5
grid. The goal is, given a handwritten image in input, to recognize the figure
this image is supposed to represent; mathematically, we look for a function
$f : [0,1]^{784} \to \{0, \ldots, 9\}$.

\begin{figure}[h]\centering
% FIGURE PLACEHOLDER (uncomment & replace with \includegraphics when image is available):
% \fbox{\parbox{0.7\linewidth}{\centering\textit{[Figure from PDF page 7: a 3x5 grid of sample MNIST handwritten digits with labels above each image (5, 0, 4, 1, 9, 2, 1, 3, 1, 4, 3, 5, 3, 6, 1).]}}}
\caption{Sample images from the MNIST dataset. Each image represents a handwritten digit.}
\label{fig:ch1-mnist}
\end{figure}

Applications can range from postal code recognition to note taking or optical
character recognition of documents.

\section{Unsupervised learning example: clustering}
\label{sec:ch1-clustering}

Clustering is an unsupervised learning technique that groups data points into
clusters based on their similarity. The data is in a form of objects $X_i$, $i \le N$.
The number $n_c$ of clusters can be specified by the user or proposed by the
algorithm. The output is thus possibly a number of clusters $n_c$ and a label
$Y_i \in \{1, \ldots, n_c\}$ for each dataset point $X_i$. The figure~\ref{fig:ch1-clustering}, adapted from scikit-learn,
illustrates the performance of various clustering algorithms on datasets
with different characteristics; all inputs are vectors $X_i \in \R^2$.

\begin{figure}[h]\centering
% FIGURE PLACEHOLDER (uncomment & replace with \includegraphics when image is available):
% \fbox{\parbox{0.7\linewidth}{\centering\textit{[Figure from PDF page 8: scikit-learn comparison grid of clustering algorithms (MiniBatch KMeans, Affinity Propagation, MeanShift, Spectral Clustering, Ward, Agglomerative Clustering, DBSCAN, HDBSCAN, OPTICS, BIRCH, Gaussian Mixture) applied to several 2D datasets (rings, moons, blobs, anisotropic, varied-variance, uniform). Labels are shown as colors.]}}}
\caption{Comparison of clustering algorithms on different datasets (source: \cite{turinici14,turinici17}). Each line is a different dataset, each column a different algorithm. Labels $Y_i$ are depicted as colors. For more details on clustering see Chapter~8.}
\label{fig:ch1-clustering}
\end{figure}

\section{General Procedure of Supervised Learning}
\label{sec:ch1-procedure}

We now describe the important objets and steps required in order to implement
a supervised learning procedure.

\subsection{Dataset}
\label{sec:ch1-dataset}

A crucially important object is the input data used to train the algorithm.
In supervised learning, the data source can be formalized as an ensemble of
pairs of input features $X$ and corresponding output labels $Y$. The dataset $\cD$
is represented as:
\begin{equation}\label{eq:ch1-dataset}
\cD = \{(X_i, Y_i)\}_{i=1}^{N} \quad \text{where} \quad X_i \in \cX,\; Y_i \in \cY.
\end{equation}

\begin{example}
$X$: Measurements of house features (e.g., size, number of
rooms), $Y$: House prices.
\end{example}

\begin{example}
$X$: Characteristics of patients (e.g., age, blood pressure), $Y$:
Diagnosis outcomes.
\end{example}

\subsection{Feature Engineering}
\label{sec:ch1-features}

Once the raw data is given an optional step is transforming this data to obtain
the good input format for the procedure. We already gave an example with transforming a MNIST image gray levels by mapping the value set
$\{0, 1, \ldots, 255\}$ to $\{0, 1/255, \ldots, 1\} \subset [0,1]$. In general such a procedure falls
within the denomination of \emph{feature engineering}. Feature engineering involves
transforming raw data into meaningful features to improve model performance.
Mathematically, this can be represented as a mapping $\phi : \cX \to \cX'$. Some examples include:
\begin{itemize}
  \item Polynomial features: $\phi(x) = [x, x^2, x^3, \ldots]$
  \item Normalization: $\phi_{\mu,\sigma}(x) = \frac{x-\mu}{\sigma}$
  \item One-hot encoding for categorical variables: $\phi(\text{category}) = [0, 1, 0, \ldots]$
  \item Log transformation: $\phi(x) = \log(x)$.
\end{itemize}

\subsection{Construction of the Predictor}
\label{sec:ch1-predictor}

The output of the model is a predictor mapping that allows to propose a
classification outcome for new objects not in the dataset. The predictor $\hat{f}$ is
constructed by selecting a model $\mathcal{M}$ and fitting it to the training data $\cD$. To
evaluate how well the model predicts the target values, we use a \emph{loss function},
which measures the discrepancy between the true output $y$ and the predicted
output $\hat{y}$. A smaller loss indicates better predictions. The model parameters
are estimated by minimizing a loss function $l(y, \hat{y})$ a process commonly referred
to as \emph{training} the model.

\begin{example}
\textbf{Linear regression} is a fundamental method in supervised
learning that models the relationship between input features and a continuous
output by assuming a linear dependency: $\mathcal{M}(X; \beta) = X\T \beta$.
\[
\hat{\beta} = \argmin_{\beta} \frac{1}{N} \sum_{i=1}^{N} l(Y_i, \mathcal{M}(X_i; \beta))
\]
\end{example}

\begin{example}
\textbf{Decision tree} is a non-parametric model that partitions the
feature space into regions using hierarchical rules, enabling flexible classification
or regression without assuming linearity. Here $\mathcal{M}(X; \theta)$ is a tree-based
model with parameters $\theta$.
\end{example}

\subsection{Model Evaluation}
\label{sec:ch1-evaluation}

Model evaluation assesses how well the predictor performs on unseen data.
This is done using metrics such as accuracy:
\begin{equation}\label{eq:ch1-accuracy}
\text{Accuracy} = \frac{\sum_{i=1}^{N} \ind_{\hat{y}_i = y_i}}{N},
\end{equation}
and by checking for common issues like \emph{overfitting} and \emph{underfitting}. Overfitting
occurs when the model performs well on training data but poorly on new
data, while underfitting means the model performs poorly on both training and new
data. These checks are crucial for ensuring \emph{generalization}, which refers to the
model's ability to maintain good performance on data it has not seen before.

\begin{tcolorbox}[advancedbox]
\textbf{Advanced topics 1.5.} \textit{To mitigate overfitting and underfitting, practitioners
often use techniques such as cross-validation to assess accuracy,
compare training and validation errors to detect overfitting, evaluate
performance on a separate test set, and monitor learning curves to
identify signs of underfitting.}
\end{tcolorbox}

\subsection{Model Interpretability}
\label{sec:ch1-interpretability}

Model interpretability involves understanding how the model makes predictions
and why certain outputs are produced. This is important for trust, debugging,
and compliance in many applications. Interpretability can be achieved
through various techniques. For instance, in a linear regression model, the
coefficients $\beta_j$ indicate the importance and direction of influence of feature $j$
on the prediction. For more complex models, such as decision trees or neural
networks, interpretability may involve feature importance scores, visualization
of decision paths, or tools like SHAP and LIME.

Another practical consideration is the trade-off between model complexity
and computational cost. Simpler models, when sufficient, are often preferable
for efficiency and transparency. For example, there is no need to deploy a large
language model (LLM) when a linear regression provides accurate results for
the task at hand!

\chapter[Useful theoretical objects]{Useful theoretical objects: predictors, loss functions, bias, variance}
\label{sec:ch2-intro}

This chapter introduces a set of theoretical objects that will be used throughout
the book, including predictors, loss functions, and the notions of bias and
variance. Together, these concepts form the basic framework for statistical
learning, allowing us to formalize prediction tasks, quantify errors, and analyze
the behavior of learning algorithms. In particular, the bias-variance
decomposition offers insight into the trade-offs that govern model complexity
and predictive performance.

\section{Predictors, Loss functions, Bayes predictor}
\label{sec:ch2-predictors}

Suppose that on $\cX$ and $\cY$ are defined some probabilistic setting, for instance
when $\cX = \R^d$ we take the Borel sigma algebra $\mathcal{B}(\R^d)$ and associated probability
measure. When $\cY = \{0,1\}$ the sigma algebra is $\mathcal{P}(\cY)$\footnote{Here and in the following for any discrete set $\mathcal{S}$ we denote $\mathcal{P}(\mathcal{S})$ the ensemble of its subsets.}.

\subsection{Predictor}
\label{sec:ch2-predictor}

In statistical machine learning, a \textbf{predictor} is a \textbf{measurable function} $f$
that maps an input $x \in \cX$ to an output $y \in \cY$. The goal of a predictor is to
accurately model the relationship between the input features and the output
variable.

\subsection{Empirical Predictor}
\label{sec:ch2-emppredictor}

Given a dataset $\cD = \{(X_i, Y_i)\}_{i=1}^{N}$, the \textbf{empirical predictor} $\hat{f}$ is some predictor
that depends on the on $\cD$ through some algorithm (also called ``model'').
This predictor approximates the underlying relationship based on the observed
data.

\subsection{Dataset hypothesis}
\label{sec:ch2-datasethyp}

We consider a probabilistic setting $(\Omega, \mathcal{F}, \PP)$ over $\Omega = \cX \times \cY$, for instance
$(\R^d \times \{0,1\}, \mathcal{B}(\R^d \times \{0,1\}))$\footnote{$\mathcal{B}(U)$ is the Borel sigma algebra associated to the set $U$ in some $\R^n$}. Denote
\begin{equation}\label{eq:ch2-nu}
\nu = \text{the distribution of the couple } (X, Y).
\end{equation}
The basic hypothesis that we make \textbf{in supervised learning} is that
\begin{equation}\label{eq:ch2-iid}
\cD \text{ is a i.i.d. sampling from the distribution } \nu \text{ in \eqref{eq:ch2-nu}.}
\end{equation}

\begin{tcolorbox}[advancedbox]
\textbf{Advanced topics 2.1.} \textit{This hypothesis is \textbf{not} true in reinforcement learning!}
\end{tcolorbox}

\subsection{Loss Function, expected loss, risk}
\label{sec:ch2-loss}

A loss function $l(\hat{y}, y)$ measures the discrepancy between the true output $y$
and the predicted output $\hat{y}$. The choice of loss function is crucial as it guides
the learning process.

\begin{definition}
We call loss function a mapping $\ell : \cY \times \cY \to \R$ such that:
\begin{align}
\ell(\hat{y}, y) &\ge 0,\ \forall y, \hat{y} \in \cY \label{eq:ch2-lossnonneg} \\
\ell(y, y) &= 0,\ \forall y \in \cY \label{eq:ch2-losszero}
\end{align}
\end{definition}

\begin{example}
For regression a common example is the mean squared error
(MSE)
\begin{equation}\label{eq:ch2-mse}
l_{MSE}(\hat{y}, y) = (y - \hat{y})^2.
\end{equation}
\end{example}

\begin{example}
For general classification we rather use the cross-entropy loss
(see later for formal definition) while for binary classification a suitable choice
is
\begin{equation}\label{eq:ch2-binary}
l_{binary}(\hat{y}, y) = \ind_{y \neq \hat{y}}.
\end{equation}
\end{example}

\begin{example}
Sometimes the following loss, called $L^1$ (or L1) loss is employed:
\begin{equation}\label{eq:ch2-l1}
l_{L1}(\hat{y}, y) = |\hat{y} - y|.
\end{equation}
\end{example}

When we sum the loss function over the whole dataset (empirical or not)
we recover an aggregated loss function depending on some predictor $\hat{f}$ that is
also called \textbf{expected loss} or \textbf{risk}:
\begin{enumerate}
  \item \textbf{empirical risk} (depends on dataset, on the loss function and on the predictor):
  \begin{equation}\label{eq:ch2-emprisk}
  \mathcal{R}^{\cD, l}(\hat{f}) = \frac{1}{N} \sum l(\hat{f}(X_i), Y_i).
  \end{equation}
  \item \textbf{risk} (depends on the predictor and the loss function):
  \begin{equation}\label{eq:ch2-risk}
  \mathcal{R}^{l}(\hat{f}) = \E[l(\hat{f}(X), Y)].
  \end{equation}
\end{enumerate}

\subsection{Bayes Predictor}
\label{sec:ch2-bayes}

The \textbf{Bayes predictor} $f_{\text{Bayes}}$ is the theoretical optimal predictor that minimizes
the risk over all possible predictors. It is defined through the property
\begin{equation}\label{eq:ch2-bayesdef}
\forall g : \cX \to \cY \text{ measurable} : \mathcal{R}^{l}(f_{\text{Bayes}}) \le \mathcal{R}^{l}(g).
\end{equation}

\begin{tcolorbox}[advancedbox]
\textbf{Advanced topics 2.6.} \textit{We will assume in the following that a Bayes
predictor exists. This assertion can be proved under some conditions on
the loss function (positivity, measurability), the joint distribution $\nu$ (to
ensure average is finite, to ensure rapid decay at infinity) and/or on the
spaces (for instance compactness).}
\end{tcolorbox}

\begin{tcolorbox}[warningbox]
\textbf{Warning 2.7.} \textit{The Bayes Predictor is not necessarily unique.}
\end{tcolorbox}

The Bayes predictor serves as a benchmark for evaluating the performance
of empirical predictors.

\begin{tcolorbox}[remarkbox]
\textbf{Intuition 2.8.} \textit{The error of the Bayes predictor is related to how much
the output can be inferred from the information in the input; for instance
it is large when the input is independent of the output and zero when the
output is a deterministic function of the input, see also equation~\eqref{eq:ch2-bayesreg}
below.}
\end{tcolorbox}

\begin{proposition}\label{prop:ch2-bayes}
Assume $\cX = \R^d$ and $\E[Y^2] < \infty$.
\begin{enumerate}
  \item For a regression problem with $l = l_{MSE}$, $Y = \R$ the Bayes predictor is:
  \begin{equation}\label{eq:ch2-bayesreg}
  f_{Br}(x) = \E[Y \mid X = x],
  \end{equation}
  where $\E[Y \mid X = x]$ is the conditional expectation of $Y$ given $x$.

  \item For a binary classification (that is $\cY = \{0,1\}$) and loss $l_{binary}$ the Bayes
  predictor is:
  \begin{equation}\label{eq:ch2-bayesbin}
  \begin{aligned}
  f_{Bc}(x) &= \begin{cases} 1, & \text{if } \PP[Y = 1 \mid X = x] \ge 1/2 \\ 0, & \text{otherwise} \end{cases} \\
  &= \sum_{y \in \{0,1\}} y \cdot \ind_{\PP[Y=y \mid X=x] \ge \PP[Y = 1-y \mid X=x]}.
  \end{aligned}
  \end{equation}
\end{enumerate}
\end{proposition}

\begin{proof}
For item 1 note that for any predictor $\hat{f}(X)$:
\begin{equation}\label{eq:ch2-decomp}
\begin{aligned}
\mathcal{R}^{l_{MSE}}(\hat{f}) &= \E[(Y - \hat{f}(X))^2] = \E\!\left[\left\{\hat{f}(X) - \E[Y \mid X]\right\}^2\right] \\
&\quad + \E[(Y - \E[Y \mid X])^2].
\end{aligned}
\end{equation}
This is true because we can decompose $Y - \hat{f}(X)$ as $Y - \E[Y \mid X] + \E[Y \mid X] - \hat{f}(X)$;
moreover the cross term is null:
\begin{equation}\label{eq:ch2-crossnull}
\E\!\left[\left\{\hat{f}(X) - \E[Y \mid X]\right\} \cdot \left\{Y - \E[Y \mid X]\right\}\right] = 0
\end{equation}
by the properties of the conditional expectation and because $\hat{f}(X) - \E[Y \mid X]$
is measurable with respect to the sigma-algebra generated by $X$. We conclude
by observing that from \eqref{eq:ch2-decomp} the risk of any predictor $\hat{f}(X)$ is larger
than that of $\E[Y \mid X]$.

For item 2 the previous relation is not useful because the predictor is constrained
to only take values in the discrete set $\cY$. Let $\hat{f}$ be some given
predictor and denote
\begin{equation}\label{eq:ch2-eta}
\eta(x) = \PP[Y = 1 \mid X = x].
\end{equation}
The error of $\hat{f}$ is
\begin{equation}\label{eq:ch2-binerror}
\begin{aligned}
R^{l_{binary}}(\hat{f}) &= \E[l_{binary}(\hat{f}(X), Y)] = \E[\ind_{Y \neq \hat{f}(X)}] \\
&= \E\!\left[\sum_{y \in \{0,1\}} \PP[Y = y, \hat{f}(X) = 1 - y]\right] \\
&= \int_{\cX} \sum_{y \in \{0,1\}} \PP[Y = y, \hat{f}(X) = 1 - y \mid X = x] p_X(dx) \\
&= \int_{\cX} \sum_{y \in \{0,1\}} \ind_{\hat{f}(x) = 1 - y} \PP[Y = y \mid X = x] p_X(dx) \\
&= \int_{\cX} \ind_{\hat{f}(x) = 1}(1 - \eta(x)) p_X(dx) + \ind_{\hat{f}(x) = 0}\eta(x) p_X(dx) \\
&\ge \int_{\cX} \min\{\eta(x), 1 - \eta(x)\} p_X(dx)
\end{aligned}
\end{equation}
with equality for the predictor which is such that $\hat{f}(x) = 0$ when $\eta(x) \le
1 - \eta(x)$ and $1$ otherwise. But this is precisely the Bayes predictor $f_{Bc}$ in
\eqref{eq:ch2-bayesbin}. In particular we proved that
\begin{equation}\label{eq:ch2-bayesbinrisk}
R^{l_{binary}}(f_{Bc}) = \E[\min\{\eta(X), 1 - \eta(X)\}].
\end{equation}
\end{proof}

Although useful, the Bayes predictor is implicit and in general not known.
One looks for algorithms / methods / procedures to approximate it.

\section{Bias-Variance Tradeoff for regression}
\label{sec:ch2-biasvar}

Consider a general regression problem where the goal is to find a predictor
$\hat{f} : \cX \to \cY$ that minimizes the expected mean square error loss (see \eqref{eq:ch2-mse}
for the definition of the loss $l_{MSE}$). Recall that in this context the \textbf{Bayes
predictor} $f_{\text{Bayes}}$ is $f_{Br}$ in \eqref{eq:ch2-bayesreg}. Note also that in this section we are not
attached to a particular dataset.

\subsection{Decomposition of the risk}
\label{sec:ch2-riskdecomp}

\begin{proposition}\label{prop:ch2-riskdecomp}
For a general predictor $\hat{f}$, the risk associated to $l_{MSE}$ can
be decomposed into two components as follows:
\begin{equation}\label{eq:ch2-riskdecomp}
\begin{aligned}
\mathcal{R}^{l_{MSE}}(\hat{f}) &= \E[(Y - \hat{f}(X))^2] \\
&= \E\!\left[\left\{\hat{f}(X) - f_{\text{Bayes}}(X)\right\}^2\right] + \E[(Y - f_{\text{Bayes}}(X))^2].
\end{aligned}
\end{equation}
\end{proposition}

\begin{proof}
This is a consequence of \eqref{eq:ch2-decomp} because in this case $f_{\text{Bayes}}(X) = \E[Y \mid X]$.
\end{proof}

\begin{tcolorbox}[remarkbox]
\textbf{Intuition 2.11.} \textit{Thus the error is composed of two terms:\\
-- a term $\E[\hat{f}(X) - f_{\text{Bayes}}(X)]^2$ that expresses the lack of optimality of
our estimator $\hat{f}$ with respect to the Bayes (optimal) estimator;\\
-- an irreducible term $\E[(Y - f_{\text{Bayes}}(X))^2]$ expressing intrinsic uncertainty
still remaining in the data and saying that information present in the
features $X$ is not enough to infer $Y$ exactly. This can only be made
lower by a change in the design variable $X$.\\
The proposition is a ``no free lunch'' type result: we can improve each
of the terms but this is often as the expense of the other.}
\end{tcolorbox}

\subsection{\texorpdfstring{Pointwise risk for dataset dependent estimators when $Y = g(X) + \varepsilon$}{Pointwise risk for dataset dependent estimators when Y = g(X) + epsilon}}
\label{sec:ch2-pointwise}

Consider the special case where the response variable $Y$ is generated as a function
of the input features plus an error term, $Y = g(X) + \varepsilon$, with $\varepsilon$ independent
of $X$ and $\E[\varepsilon] = 0$, $\Var[\varepsilon] = \sigma^2$. In this case, the Bayes predictor simplifies to:
$f_{\text{Bayes}} = g$.

We take a particular type of estimator $\hat{f}^{\cD}$, namely one that is produced by a
procedure operating on the dataset $\cD$ with input features $X_i \in \cX$ and response
variables $Y_i \in \cY$. We honor hypothesis in Section~\ref{sec:ch2-datasethyp} and consider $X_i, Y_i$ as
i.i.d. realization of some random couple $(X, Y)$. \textbf{It is the uncertainty in
this realization that is considered in the following} i.e., we recognize
that for a given, \textbf{fixed} $x_v \in \cX$ the stochasticity in $X_i, Y_i$ thus in $\cD$ will induce
a stochasticity in $\hat{f}^{\cD}(x_v)$.

\begin{proposition}\label{prop:ch2-biasvar}
Under the previous assumptions (in particular the loss
being $l_{MSE}$, $Y = g(X) + \varepsilon$, \ldots) the bias-variance decomposition for the predictor
$\hat{f}^{\cD}$ at some point $x_v$ is:
\begin{equation}\label{eq:ch2-biasvardecomp}
\E[(Y - \hat{f}^{\cD}(x_v))^2] = \underbrace{(g(x_v) - \E[\hat{f}^{\cD}(x_v)])^2}_{(\mathrm{Bias}[\hat{f}^{\cD}(x_v)])^2} + \underbrace{\E[(\hat{f}^{\cD}(x_v) - \E[\hat{f}^{\cD}(x_v)])^2]}_{\Var[\hat{f}^{\cD}(x_v)]} + \sigma^2.
\end{equation}
\end{proposition}

\begin{remark}
Here $x_v$ and $g(x_v)$ are fixed, deterministic while $Y_v$ is random
because comes from the relation $Y = g(x_v) + \varepsilon$ and $\hat{f}^{\cD}$ is also random
because $\cD$ is random. Also important is to note that $Y$ and $\cD$ are \textbf{independent}.
\end{remark}

\begin{proof}
We follow several steps:

\textbf{Step 1}: Start with the definition of $Y$: $Y = g(x_v) + \varepsilon$ and write the expected
squared error
\[
\E[(Y - \hat{f}^{\cD}(x_v))^2] = \E[(g(x_v) + \varepsilon - \hat{f}^{\cD}(x_v))^2].
\]

\textbf{Step 2}: Expanding the square and using the properties of expectation we
obtain
\[
\E[(g(x_v) - \hat{f}^{\cD}(x_v))^2] + \E[\varepsilon^2] + 2\E[\varepsilon(g(x_v) - \hat{f}^{\cD}(x_v))].
\]

\textbf{Step 3}: Take as previously a conditional expectation with respect to $\cD$ and
since on one hand $\varepsilon$ is independent of $\cD$ and on the other hand $\E[\varepsilon] = 0$,
the cross-term vanishes; we are left with $\E[(g(x_v) - \hat{f}^{\cD}(x_v))^2] + \sigma^2$.

\textbf{Step 4}: Recall the usual bias-variance decomposition true for any random
variable $Z$ and any $a \in \R$:
\begin{equation}\label{eq:ch2-bvgeneric}
\E[(Z - a)^2] = \underbrace{(\E[Z] - a)^2}_{\text{bias squared}} + \underbrace{\E[(Z - \E[Z])^2]}_{\text{variance } \Var[Z]}.
\end{equation}

We set $Z = \hat{f}^{\cD}(x_v)$ and $a = g(x_v)$ to obtain:
\begin{equation}\label{eq:ch2-biasvarfinal}
\E[(g(x_v) - \hat{f}^{\cD}(x_v))^2] = (g(x_v) - \E[\hat{f}^{\cD}(x_v)])^2 + \E[(\hat{f}^{\cD}(x_v) - \E[\hat{f}^{\cD}(x_v)])^2].
\end{equation}
Here we identify the first part as $\left(\mathrm{Bias}[\hat{f}^{\cD}(x_v)]\right)^2$ and the second as
$\Var[\hat{f}^{\cD}(x_v)]$.
\end{proof}

\begin{figure}[h]\centering
% FIGURE PLACEHOLDER (uncomment & replace with \includegraphics when image is available):
% \fbox{\parbox{0.7\linewidth}{\centering\textit{[Figure from PDF page 21: four-quadrant illustration of the bias-variance decomposition. Rows labeled ``Low Bias'' / ``High Bias'', columns labeled ``Low Variance'' / ``High Variance''. Each quadrant shows a target (concentric circles) with dart-like marks illustrating the distribution of predictions relative to the center.]}}}
\caption{Illustration of the bias-variance decomposition in Proposition~\ref{prop:ch2-biasvar}.}
\label{fig:ch2-biasvar}
\end{figure}

\begin{tcolorbox}[advancedbox]
\textbf{Advanced topics 2.14.} Practical implications of bias-variance
decomposition \eqref{eq:ch2-biasvardecomp} (see also Figure~\ref{fig:ch2-biasvar}):
\begin{itemize}
  \item the term $(g(x_v) - \E[\hat{f}^{\cD}(x_v)])^2$ represents the \textbf{bias squared}, indicating
  the error due to the simplifications in the model. It is called ``bias''
  to convey our hope that somehow the average over all datasets of our
  model $\hat{f}^{\cD}$ will be $f_{\text{Bayes}}$. If this is not true then our procedure is biased.
  High bias can lead to \textbf{underfitting}.
  \item the term $\E[(\hat{f}^{\cD}(x_v) - \E[\hat{f}^{\cD}(x_v)])^2]$ represents the \textbf{variance} of the predictor
  $\hat{f}^{\cD}$. It says that when the dataset changes the predictor will also
  change so it reflects the dependence of the predictor on the particular
  choice of dataset $\cD$ rather than on its general statistical properties. Often,
  when we want to reduce the bias the model learns too much of the
  details of the dataset $\cD$ and this variance will be high which can lead to
  \textbf{overfitting} (the model is too sensitive to the training data).
  \item finally, the term $\sigma^2$ represents the \textbf{irreducible error} due to the inherent
  noise in the data. This term sets a lower bound on the achievable
  error.
\end{itemize}
In practice, regularization techniques such as ridge regression and
LASSO are employed to manage the bias-variance tradeoff. They add
a regularization term to the loss function; the coefficient of this term is
used to control the bias-variance trade-off.
\end{tcolorbox}

\section{Exercises}
\label{sec:ch2-exercises}

\begin{tcolorbox}[remarkbox]
\textbf{Exercise 2.1} (Bayes predictor \& mensurability). \textit{Show that if $Y$ is $X$
measurable then the error of the Bayes predictor is null.}
\end{tcolorbox}

\begin{tcolorbox}[remarkbox]
\textbf{Exercise 2.2} ($L^1$ Bayes predictor). \textit{Consider the $L1$ loss \eqref{eq:ch2-l1}.}
\begin{enumerate}
  \item \textit{When $\cY = \R$ find the Bayes predictor for $l_{L1}$.}
  \item \textit{What happens when $\cY = \R^m$ for some integer $m \ge 2$?}
\end{enumerate}
\end{tcolorbox}

\begin{tcolorbox}[remarkbox]
\textbf{Exercise 2.3} (Bayes predictor \& independence). \textit{Let $\ell$ denote a general
loss function and suppose that $X$ and $Y$ are independent. Define the
function $g : \cY \to \R$ by:}
\begin{equation}\label{eq:ch2-exindep}
g(a) := \E_Y[\ell(a, Y)],
\end{equation}
\textit{and suppose $g$ admits a global minimum $a_{min}$ on $\cY$.}
\begin{enumerate}
  \item \textit{Show that $f(x) = a_{min}\ \forall x \in \cX$ is a Bayes predictor.}
  \item \textit{Explain how the result in item 1 applies for $\ell = l_{MSE}$ defined in \eqref{eq:ch2-mse}.}
  \item \textit{Explain how the result in item 1 applies when $Y$ is Bernoulli with
  $\PP[Y = 1] = p = 1 - \PP[Y = 0]$ and $\ell = l_{binary}$ defined in \eqref{eq:ch2-binary}.}
\end{enumerate}
\end{tcolorbox}
